% ===========================================================================
%  Epidemic Network Spread -- KDD Explorations (double column) draft
%  Class: sigkddExp.cls   Bibliography: references.bib
%  Build: tectonic main.tex   (runs bibtex automatically)
%  STATUS: pre-results working draft. % TODO marks need data/figures.
% ===========================================================================
\documentclass{sigkddExp}

\begin{document}

\title{Topology-Aware Immunization of the European Air-Travel Network:
Compartmental Models and the Degree--Betweenness Question}

\numberofauthors{5}
\author{
\alignauthor Author One\\
       \affaddr{MSc Network Science, Your University}\\
       \email{author.one@example.edu}
\alignauthor Author Two\\
       \affaddr{MSc Network Science, Your University}\\
       \email{author.two@example.edu}
\alignauthor Author Three\\
       \affaddr{MSc Network Science, Your University}\\
       \email{author.three@example.edu}
}
\additionalauthors{Additional authors: Author Four and Author Five
(Your University, \texttt{authors@example.edu}).}
\date{\today}
\maketitle

\begin{abstract}
The topology of a mobility network decides how an epidemic spreads and
which intervention contains it. We model disease propagation on the
European subgraph of the OpenFlights air-travel network as a
metapopulation reaction--diffusion process, running four compartmental
dynamics (SIR, SIS, SEIR, SQIR) with parameters anchored to literature
reproduction numbers, and we compare vaccination by random, degree, and
betweenness targeting. Our central, testable question is where Europe sits
on a known spectrum: in the US domestic network the largest hubs are also
the main bridges, so degree- and betweenness-targeted immunization
coincide, whereas in the worldwide network geopolitical community
structure produces anomalous high-betweenness, low-degree gateways that
only betweenness targeting captures. A subgraph- and core-aware layer
($k$-cores, motifs, graphlets) refines targeting and explains outcomes.
The study is a reproducible pipeline whose region and transport layers
(air, land, water) are parameters, so the same analysis extends across
continents and to multimodal combinations. Claims are comparative, not
predictive.
\end{abstract}

\section{Introduction}
A single immunization policy cannot be optimal for every pathogen. We ask
how vaccination strategy interacts with both the infection model and the
topology of the travel network, on a real, strongly heterogeneous
substrate. Concretely, we simulate spread along European aviation routes
and test random, degree-, and betweenness-targeted vaccination across four
compartmental models, reinforcing the principle that a strategy should be
matched to the available infection information rather than applied
uniformly. Modelling epidemics on flight networks is not itself new; our
specific positioning and the gap we fill are stated in
Section~\ref{sec:novelty}.

\section{Related Work}
Compartmental epidemics originate with Kermack and
McKendrick~\cite{kermack:sir}; network epidemiology asks how contact
structure changes those dynamics~\cite{pastorsatorras:review,keeling:networks}.
Scale-free topology can destroy the epidemic
threshold~\cite{pastorsatorras:scalefree}, and Newman linked SIR to bond
percolation for exact outbreak results~\cite{newman:spread}. Because hubs
dominate spreading, targeted immunization far outperforms
random~\cite{pastorsatorras:immunization,cohen:immunization}, though the
best spreaders sit in the network \emph{core} rather than simply at the
highest degree~\cite{kitsak:spreaders}. At the inter-city scale, the air
network governs global spread~\cite{colizza:airline,balcan:multiscale},
formalised as a metapopulation reaction--diffusion process with a global
invasion threshold~\cite{colizza:reactiondiffusion,colizza:invasion}.
Structurally, motifs~\cite{milo:motifs} and graphlet
signatures~\cite{przulj:graphlets} fingerprint local topology but are
rarely used to steer immunization --- the gap we target.

\section{Novelty and Positioning}
\label{sec:novelty}
\textbf{(1) Europe, four models, one protocol.} Prior work targets the
worldwide or US networks; a Europe-focused comparison across SIR, SIS,
SEIR and SQIR under a common protocol is, to our knowledge, missing.

\textbf{(2) The degree--betweenness question.} In hub-and-bridge networks
such as the US domestic system, degree and betweenness are highly
correlated and the two immunization strategies give nearly identical
results~\cite{tanaka:centrality}. % TODO: verify Tanaka 2014 reference.
On the worldwide network, Guimer{\`a} et al.~\cite{guimera:anomalous} found
\emph{anomalous} centrality: community and geopolitical structure give some
low-degree airports very high betweenness (Anchorage being canonical).
Where a network sits on this spectrum decides whether the strategies
converge or diverge. \emph{We test empirically where Europe lies.} Because
region is a pipeline parameter (Section~\ref{sec:pipeline}), we place
Europe on the spectrum directly by repeating the analysis across continents
and comparing the degree--betweenness correlation.

\textbf{(3) Subgraph- and core-aware layer.} We use $k$-core decomposition
and small-subgraph statistics both to target immunization and to explain
why a model--strategy pair behaves as it does, connecting the result to the
local structure that produces it.

\section{Data and Network}
\label{sec:data}
The primary substrate is the air layer, from the OpenFlights route
database~\cite{openflights}, built with NetworkX~\cite{hagberg:networkx}:
nodes are airports, edges are routes weighted by flight frequency. The
region is a parameter of network generation (Section~\ref{sec:pipeline}),
so the European subgraph is the focus while other continents and the world
follow for the cross-region comparison.

\textbf{Multimodal data sources.} The same node set admits ground and sea
layers. For \emph{land}, rail and road \emph{topology} is globally available from
OpenStreetMap and the GRIP roads database~\cite{meijer:grip}. To keep the
cross-region comparison valid, \emph{flow} weights must be built the
\emph{same way for every region}; we therefore generate them with a single
gravity/radiation model calibrated on population~\cite{simini:radiation},
applied uniformly, and use Eurostat commuting data only to \emph{validate}
that model within Europe --- never as the production source for one region
while others differ. Commuting dominates short-range flux and the proxy
choice materially affects predictions~\cite{balcan:multiscale,tizzoni:proxies}.
For \emph{water}, we use OpenStreetMap ferry routes and open cargo
port-call data, whose network structure and disease relevance are
documented in~\cite{kaluza:cargo,seebens:bioinvasion,mari:cholera}. Air is
in use; land and water enter as additional tagged layers.

\textbf{Modelling assumptions.} (i) Population scales with degree,
\begin{equation}
\mathrm{pop}(v) = P_{0} + d(v)\,P_{\text{route}},
\end{equation}
with $P_{0}=150{,}000$ and $P_{\text{route}}=45{,}000$, assuming more routes
imply a larger city --- a proxy, not census ground truth. (ii) Interventions
are biological shields only: vaccinated cities become immune dead-ends and
no nodes or edges are removed, isolating the effect of vaccination from
route closures. (iii) The travel rate $\tau$ is scaled by route frequency,
so busier corridors carry proportionally more travellers. (iv) The outbreak
starts from a single random patient-zero seed, averaged over seeds. (v) The
network is static (no timetables or seasonality). The resulting European
grid is a hybridised scale-free network with extreme degree heterogeneity,
in which hubs (London, Paris, Frankfurt, Amsterdam) connect high-volume
corridors.
% TODO: report measured N, E, <k>, <k^2>/<k>, assortativity.

\section{Models and Parameters}
\subsection{Metapopulation reaction--diffusion}
Each airport carries a full compartmental model and individuals diffuse
along edges~\cite{colizza:reactiondiffusion}. Per day, a \emph{reaction}
step applies the force of infection $\lambda_i=\beta\,I_i/N_i$ with
model-specific transitions at rates $\gamma$ (recovery), $\sigma$
(incubation, SEIR) and $\kappa$ (quarantine, SQIR); a \emph{diffusion} step
moves a fraction of each compartment along outbound routes, split by route
weight $w_{ij}$ and a global travel rate $\tau$. The four models span
permanent immunity (SIR: measles), no immunity (SIS: common cold), latency
(SEIR: COVID-19, SARS) and isolation (SQIR: Ebola). This formulation
carries a global invasion threshold below which spread stays
local~\cite{colizza:invasion}, and is substrate-agnostic, so it applies
unchanged to land and water layers.

\subsection{Parameter choice}
We fix clinical rates ($\gamma,\sigma,\kappa$) from each disease's
epidemiology and set $\beta$ to a literature $R_0$. On a heterogeneous
network the well-mixed identity $R_0=\beta/\gamma$ fails: degree
fluctuations inflate it as
$R_0\propto\frac{\beta}{\gamma}\frac{\langle k^2\rangle}{\langle k\rangle}$
\cite{pastorsatorras:scalefree}, and for multi-compartment models we
compute $R_0$ via the next-generation matrix~\cite{vandendriessche:r0}.

\begin{table}[t]
\centering
\caption{Per-model parameters used in the experiments. Rates are per day;
$R_0^{\text{loc}}=\beta/\gamma$ is the well-mixed local reproduction number
(the network value is larger, scaled by $\langle k^2\rangle/\langle k\rangle$).}
\label{tab:params}
\begin{tabular}{|l|c|c|c|c|}\hline
 & SIR & SIS & SEIR & SQIR\\\hline
Transmission $\beta$ & 0.32 & 0.32 & 0.32 & 0.32\\\hline
Recovery $\gamma$ & 0.12 & 0.12 & 0.12 & 0.12\\\hline
Incubation $\sigma$ & --- & --- & 0.20 & ---\\\hline
Quarantine $\kappa$ & --- & --- & --- & 0.15\\\hline
$R_0^{\text{loc}}$ & 2.7 & 2.7 & 2.7 & 2.7\\\hline
\end{tabular}
\end{table}

\subsection{Calibration vs.\ verification}
We do \emph{not} calibrate $R_0$ to case data as predictive frameworks
do~\cite{balcan:gleam}, because the OpenFlights substrate carries no
ground-truth outbreak to fit. Instead we (i) select an operating point in
the $R_0$ regime where strategies actually differ, relative to the
threshold~\cite{colizza:invasion}; (ii) run a sensitivity sweep showing the
\emph{ranking} of strategies is stable across $\beta$, $\tau$, coverage and
efficacy; and (iii) verify the implementation against analytic limits
(single-node SIR; bond-percolation final size~\cite{newman:spread}). This
recovers calibration's credibility without fitting, and keeps our claims
comparative rather than predictive of absolute case counts.

\section{Vaccination Strategies}
With a fixed budget of 15 cities at $75\%$ coverage and $85\%$ efficacy,
we compare: \textbf{control} (none); \textbf{random}; \textbf{degree}
(densest traffic); and \textbf{betweenness} (structural bridges). We also
contrast low ($1\%$) versus high ($75\%$) coverage. As a structural
refinement we add a \textbf{subgraph/core} strategy: prioritise innermost
$k$-shells~\cite{kitsak:spreaders} and cities embedded in dense
transmission motifs~\cite{milo:motifs,przulj:graphlets} rather than
isolated high-degree nodes; the same statistics explain why an
intervention dismantles the core or merely trims its periphery.

\textbf{Expected mechanisms (hypotheses).} The same strategy should act
differently per model, which is the point of testing all four. Targeting
high-betweenness hubs severs Europe's structural
bridges~\cite{pastorsatorras:immunization}, and we hypothesise: for
\emph{SIR} (permanent immunity) this boxes the outbreak into its origin
region, where it starves as local susceptibles deplete --- a firebreak; for
\emph{SIS} (no immunity), which cannot be eradicated and reaches an endemic
plateau, hub immunization acts instead as a continuous filter that
compresses the active baseline; for \emph{SEIR} (latency), where
pre-symptomatic passengers defeat departure screening, pre-vaccinating hubs
intercepts them at layover nodes; and for \emph{SQIR} (isolation), shifting
identified hub cases into quarantine cuts the busiest corridors. These are
hypotheses to be tested, not results.

\section{Pipeline and Substrate}
\label{sec:pipeline}
The study is a reproducible three-stage pipeline --- data retrieval,
network generation, evaluation --- in which a run is a pure function of its
configuration and seed, and the sweep (regions $\times$ layer combinations
$\times$ models $\times$ strategies $\times$ coverage $\times$ seeds) is
orchestrated with Nextflow~\cite{ditommaso:nextflow}. Because the reaction
step is substrate-agnostic, network generation emits every combination of
the air, land and water layers of Section~\ref{sec:data}, making multimodality a
comparable axis. Layers are kept tagged with their own travel rates, since
coupling can cross the epidemic threshold even when no single layer
would~\cite{colizza:invasion}.

\section{Results}
\label{sec:results}
We build the European air network (561 airports, 10{,}072 directed routes;
$\langle k\rangle = 35.9$, $\langle k^2\rangle/\langle k\rangle = 117$,
confirming strong degree heterogeneity). The simulation horizon is set to
$210$ days from the operating-point scan rather than by convention: at
$\tau=0.0002$ the latest-peaking model (SEIR) peaks at day $\approx 169$
and SIR at day $\approx 90$, so shorter horizons truncate every model
before its peak.

\textbf{The degree--betweenness question.} Europe sits at the
\emph{correlated} end of the centrality spectrum, with Spearman
$\rho(\text{degree},\text{betweenness}) = 0.90$ and only $9$ anomalous
low-degree/high-betweenness gateways (e.g.\ Kittil\"a, Isles of Scilly,
Ivalo). Running the identical construction across regions confirms a true
spectrum (Table~\ref{tab:regions}): Europe and Asia are the most
US-like, while the Americas ($\rho=0.73$, $59$ gateways) and Oceania
($\rho=0.67$) are the most worldwide-like (anomalous). Europe is therefore
correlated but not extreme --- it carries its own peripheral gateways.

\begin{table}[t]
\centering
\caption{Cross-region degree--betweenness correlation (air layer).}
\label{tab:regions}
\begin{tabular}{|l|c|c|}\hline
Region & $\rho(\text{deg},\text{btw})$ & anomalous gateways\\\hline
Europe & 0.90 & 9\\\hline
Asia & 0.85 & 16\\\hline
Africa & 0.78 & 7\\\hline
Americas & 0.73 & 59\\\hline
Oceania & 0.67 & 16\\\hline
\end{tabular}
\end{table}

\textbf{Vaccination strategies.} Across all four models, targeted
allocation outperforms random, which barely beats no intervention. For SIR
at high coverage the mean reduction in peak active infections relative to
control is: degree $-11.4\%$, subgraph $-11.2\%$, betweenness $-9.0\%$,
$k$-core $-8.4\%$, random $-1.4\%$. That degree edges out betweenness is
consistent with Europe's high $\rho$: where the two centralities largely
agree, both target the same hubs, so the choice matters less --- precisely
the regime-dependence the paper set out to test.

\section{Conclusions}
We posed a Europe-specific, multi-model test of whether immunization
strategy must be matched to network topology, centred on the
degree--betweenness question. Empirically Europe is degree--betweenness
correlated ($\rho=0.90$), so degree- and betweenness-targeting nearly
coincide and both beat random allocation by roughly $9$--$11\%$ in peak
suppression; the divergence the literature reports for the worldwide
network appears, in our cross-region comparison, in the Americas and
Oceania rather than in Europe. \emph{Limitation:} populations and the
travel rate are synthetic proxies, so these are comparative, not
predictive, results; the land/water layers currently use radiation-model
flows rather than measured rail/road geometry.

\bibliographystyle{abbrv}
\bibliography{references}

\end{document}
